\emph{(1).} By \claimref{synthesis.rung-field-d-term-action-and-tree-masses-on-the-lift-locus}(1) the $U(2)$-invariant polynomials are generated by $|\alpha|^{2}$, $|\beta|^{2}$, $Q$ and $\bar{Q}$, each of degree two in $c$, so the invariant polynomials homogeneous of degree four are spanned by the products of two generators; their real parts and imaginary parts give the ten $I_{i}$ (the products $|\alpha|^{2}\bar{Q}$, $\bar{Q}^{2}$ and so on add nothing over $\R$). On $L$, $Q = -H^{2}$, $|\alpha|^{2} = \mathsf{N} + m$ and $|\beta|^{2} = \mathsf{N} - m$, so $I_{1} = (\mathsf{N} + m)^{2}$, $I_{2} = (\mathsf{N} - m)^{2}$, $I_{3} = \mathsf{N}^{2} - m^{2}$, $I_{4} = -H^{2}(\mathsf{N} + m)$, $I_{6} = -H^{2}(\mathsf{N} - m)$, $I_{8} = I_{10} = H^{4}$ and $I_{5} = I_{7} = I_{9} = 0$; collecting terms gives the stated map $k \mapsto (a, b, e, d, f, c_{0})$. The conjugation $c^{*} = (\bar{\beta}, \bar{\alpha})$ exchanges $|\alpha|^{2}$ and $|\beta|^{2}$ and conjugates $Q$, so it exchanges $k_{1} \leftrightarrow k_{2}$, $k_{4} \leftrightarrow k_{6}$ on $L$, which reverses $e$ and $f$ and fixes the rest. The D-term $V_{D} = \tfrac12\bigl(g^{2}\sum_{j}\mu_{j}^{2} + g'^{2}(\mu_{Y} - \zeta)^{2}\bigr)$ has the terms $-g'^{2}\zeta\mu_{Y}$ of degree two and $\tfrac12 g'^{2}\zeta^{2}$ of degree zero, so it is homogeneous of degree four exactly when $\zeta = 0$; then on $L$ it is $\tfrac18 g^{2}(\mathsf{N}^{2} - H^{4}) + \tfrac18 g'^{2}m^{2}$ (stored clause (2)). The check verifies the invariance under five rational $U(2)$ elements, the restriction at four points of $L$ for three coefficient vectors, the conjugation rule, that the six monomials $\mathsf{N}^{2}, m^{2}, \mathsf{N}m, H^{2}\mathsf{N}, H^{2}m, H^{4}$ are independent on $L$ (rank six on six points), and the homogeneity test of $V_{D}$ for $\zeta = 0$ and $\zeta = \tfrac13$.

\emph{(2).} On the family $\mathsf{N} = H^{2}\cosh t$ and $m = -H^{2}\sinh t$ (\claimref{synthesis.rung-field-on-the-lift-locus}(4)), and substituting into (1) with $\cosh^{2}t = 1 + \sinh^{2}t$ gives the family formula. The rescaling $c_{\lambda}(H) \mapsto s\,c_{\lambda}(H) = c_{s\lambda}(sH)$ is immediate from $c_{\lambda} = -\lambda w + (H^{2}/\lambda)v$, and $V(sc) = s^{4}V(c)$ by homogeneity, so $V(c_{\lambda})/H^{4}$ is a function of $\lambda/H$, that is, of $t$. Write $X = e^{t}$, $u = \cosh t$, $s = \sinh t$, $F(t) = (a + b)s^{2} - esu + du - fs$. Then $F'(t) = (a + b)\sinh 2t - e\cosh 2t + ds - fu$ and $(a + b)\sinh 2t - e\cosh 2t = \tfrac12(K_{+}X^{2} - K_{-}X^{-2})$. If $K_{+} < 0$ the leading term $\tfrac14 K_{+}X^{2}$ of $F$ wins over the terms of order $X$ and $X^{-2}$, so $F \to -\infty$. At a critical point with $t \ge 0$, $K_{+}X^{2} = K_{-}X^{-2} + 2(fu - ds) \le |K_{-}| + 2(|d| + |f|)X$, since $u, s \le X$ and $X^{-2} \le 1$; dividing by $X^{2}$ gives the first inequality. For $K_{+} > 0$ put $D = |d| + |f|$, $R = \sqrt{K_{+}|K_{-}|}$ and $X_{0} = (2D + R)/K_{+}$: then $K_{+}X_{0}^{2} - 2DX_{0} - |K_{-}| = X_{0}R - R^{2}/K_{+} = 2DR/K_{+} \ge 0$, and $K_{+}X^{2} - 2DX$ is increasing for $X > X_{0}$, so $2F'(t) \ge K_{+}X^{2} - 2DX - |K_{-}| > 0$ there: no critical point lies above $X_{0}$. With $\cosh t \le e^{t}$, $e^{-t} \le H^{2}/\mathsf{N}^{*}$, which gives the stated condition at the weak scale; the observed weak scale gives $\mathsf{N}^{*}/H^{2} = \tfrac12(v/H)^{2} \approx 2 \times 10^{88}$ with $H^{2} = \Lambda/3$ ($v/H \approx 2.1 \times 10^{44}$), or about $1.4 \times 10^{88}$ if $H$ is read as $H_{0}$ ($v/H_{0} \approx 1.7 \times 10^{44}$, \claimref{synthesis.rung-unit-second-free-number}, not used in the proof); the two readings differ by $\sqrt{\Omega_{\Lambda}}$, a factor of order one. The bound is attained up to a factor two: with $d < 0 < f$ and $f > |d|$ criticality gives $K_{+} = D/X + (f - |d|)X^{-3} + K_{-}X^{-4} > D/X$, so no bound of the form $K_{+} \le D/X$ holds. The half $t \le 0$ follows by applying the conjugation, which maps $t \mapsto -t$ (\claimref{synthesis.rung-field-on-the-lift-locus}(5)) and $(e, f) \mapsto (-e, -f)$. The check verifies the family formula at four frames for three coefficient vectors, the rescaling, the identity for the leading part of $F'$, $F'$ against a difference quotient of $F$, positivity of $F'$ on forty points above a rational upper bound of $X_{0}$ for three samples, the inequality at three exactly constructed critical points (with $f$ solved for criticality, one with $K_{+} = 10^{-3}$ and $X = 10^{3}$), two near-saturating critical points (with $K_{+}$ solved for criticality) at which $K_{+}$ exceeds $D/X$, and the runaway for two samples with $K_{+} < 0$.

\emph{(3).} As $H \to 0$ the terms $H^{2}(d\mathsf{N} + fm)$ and $c_{0}H^{4}$ drop, the locus becomes the cone $Q = 0$ and the orbit space $\mathsf{N}^{2} - m^{2} \ge H^{4}$ becomes $|m| \le \mathsf{N}$. Write $m = -\sigma \mathsf{N}$, $\sigma \in [-1, 1]$; then $V_{4} = \mathsf{N}^{2}\phi(\sigma)$ with $\phi(\sigma) = a + b\sigma^{2} - e\sigma$, the physical ray is $\sigma = 1$ and the mirror ray $\sigma = -1$. Vanishing on the physical ray is $\phi(1) = K_{+} = 0$. Given that, $\phi(-1) = 2e$, so $e \ge 0$ is necessary; $\phi'(1) = 2b - e = e - 2a$, and if it were positive $\phi$ would be negative just below $\sigma = 1$, so $2a \ge e$ is necessary. Conversely, if $b = e - a \le 0$ then $\phi$ is concave and its minimum on $[-1, 1]$ is at an endpoint, $\min(0, 2e) = 0$; if $b > 0$ the vertex $\sigma_{0} = e/(2b)$ satisfies $\sigma_{0} \ge 1$ exactly when $e \ge 2(e - a)$, that is $2a \ge e$, so $\phi$ decreases on $[-1, 1]$ to $\phi(1) = 0$. The masses: at $c = (0, (-\lambda, 0)) \in L_{0}$ the directions $d_{n} = (0, (-1, 0))$ (the rung unit) and $d_{c} = ((0, z), 0)$, $z \in \C$, are tangent to $L_{0}$ and orthogonal for $G$ to the gauge orbit, and the complement of the gauge orbit in the tangent space has dimension three (stored construction, \claimref{synthesis.rung-field-on-the-lift-locus}(3), here on $L_{0}$). Along $d_{n}$, $V_{4} = K_{+}\mathsf{N}^{2} = 0$, so the neutral mode is massless. Along $d_{c}$ exactly, $\mathsf{N}$ and $m$ both increase by $\tfrac12\varepsilon^{2}|z|^{2}$, so $V_{4}$ increases by $\tfrac12\varepsilon^{2}|z|^{2}(\partial_{\mathsf{N}} + \partial_{m})V_{4} = \tfrac12\varepsilon^{2}|z|^{2}\cdot 2(a - b)\mathsf{N}$ at $m = -\mathsf{N}$; with the kinetic coefficient $2\mathsf{N}(d_{c}) = |z|^{2}$ this is $m_{\pm}^{2} = 2(a - b)\mathsf{N} = 2(2a - e)\mathsf{N}$, and $2\mathsf{N} = v^{2}$ by $M_{W}^{2} = \tfrac12 g^{2}\mathsf{N} = \tfrac14 g^{2}v^{2}$. The check computes these Hessians on $L_{0}$ by exact interpolation along curves kept on $Q = 0$ to second order, for three admissible coupling sets and two points of the ray, and verifies that the mixed second derivative of neutral and charged directions vanishes.

\emph{(4).} By (1), the level-zero D-term plus $\kappa \mathsf{N}^{2}$ has $a = \tfrac18 g^{2} + \kappa$, $b = \tfrac18 g'^{2}$, $e = 0$, so $K_{+} = 0$ forces $\kappa = -\tfrac18(g^{2} + g'^{2})$ and then $2a - e = -\tfrac14 g'^{2} < 0$; by (3), $m_{\pm}^{2} = -\tfrac12 g'^{2}\mathsf{N} = M_{W}^{2} - M_{Z}^{2}$, which is the stored sum rule (\claimref{synthesis.rung-field-charged-scalar-sum-rule}(ii)) at $m_{h} = 0$, and $\phi(0) = a < 0$ is a negative value inside the cone. With $\kappa'm^{2} + e\mathsf{N}m$ added, $b = \tfrac18 g'^{2} + \kappa'$; admissibility requires $b = e - a \le e - \tfrac12 e$. The check computes $V_{D}$ from the moment maps (not from the reduced form), on $L_{0}$, for $g = \tfrac23$, $g' = \tfrac12$: flatness along the ray, the charged mass squared $M_{W}^{2} - M_{Z}^{2}$ by exact Hessian, a negative value at an interior point of the cone, and with $\kappa' = -\tfrac18 g'^{2} - \tfrac15$ the positive charged mass $4a\mathsf{N}$.

\emph{(5).} This clause restates the stored $X_{\mu} = B_{1/\mu}\circ X_{1}$ (\claimref{synthesis.rung-unit-transverse-to-norm}(1)) through the standard fact, cited from memory (D. Hestenes; C. Doran and A. Lasenby, Geometric Algebra for Physicists), that a boost of the null pair is the dilator of the conformal model. $\mu X_{1}(x/\mu) = \mu n_{0}' + x + \tfrac{Q(x)}{2\mu}n_{\infty}' = X_{\mu}(x)$ by $Q(x/\mu) = Q(x)/\mu^{2}$; the map $x \mapsto X_{1}(x)$ is the conformal embedding of the ground with null pair $(n_{0}', n_{\infty}')$, and $X_{1}(x/\mu)$ is the embedding of the dilated point, projectively equal to $X_{\mu}(x)$. The rung component of $X_{\mu}(I)$ is $c_{\mu}$ (\claimref{synthesis.rung-unit-transverse-to-norm}(1)), the family of (2). The same stored boost is read as the flow of $2T_{3R}$ at real rapidity on the ten in \claimref{synthesis.majorana-scale-not-carried-by-rung-frame-boost}; on the rung plane the two readings agree. The check builds the conformal model of a ground plane in the exact Clifford engine and verifies nullity of $X_{\mu}(x)$, the identity at six pairs $(x, \mu)$, and the rung component with $Q(c_{\mu}) = -H^{2}$.

\emph{Remarks.} (a) On $L$ itself the flat condition $K_{+} = 0$ with $d = f = 0$ leaves $F(t) = -\tfrac{e}{2}(1 - e^{-2t})$, a slope of relative order $H^{4}/\mathsf{N}^{2}$ along the boost hyperbola: the constraint's scale $H^{2}$ is negligible against any quantum correction at the weak scale, but it is a dimensionful input, and the flat direction on $L$ is the hyperbola $\mathsf{N}^{2} - m^{2} = H^{4}$, which (3) treats in its limit as a ray. (b) This result does not compute the running of $K_{+}$: dimensional transmutation trades the level $\xi$ for the scale at which $K_{+}$ vanishes, which is a further input unless a boundary value of the couplings at a fixed scale is supplied.
